% (User input window)

\paragraph{整数点谱率、原子 Witt 同余分层与解析--RH 分离}
附录 \ref{app:pom-projective-pressure-operator} 已把整数 \(q\) 的 moment-kernel
\(\widetilde A_q\) 提升到实参数射影转移算子 \(\mathcal L_q\)；
附录 \ref{app:real-input-40-zeta-u} 与
\ref{app:real-input-40-finite-rh} 则把 primitive 的原子 Witt 分量、
解析半径、Dirichlet 扭曲与平方根误差窗口全部写成可核验的代数对象。
下列条目记录这些接口在当前文稿中的几条直接结构后果。

\begin{theorem}[整数 moment-kernel 谱半径满足次乘法律]\label{thm:xi-projective-moment-radius-submultiplicative}
沿用定义 \ref{def:pom-projective-transfer-operator}、
命题 \ref{prop:pom-projective-operator-invariant-homogeneous}
与定理 \ref{thm:pom-projective-operator-degenerates-to-moment-kernel} 的记号。
对每个整数 \(q\ge 2\)，记
\[
r_q:=\rho(\widetilde A_q)=\rho(\mathcal L_q|_{\mathcal H_q}).
\]
则对任意整数 \(a,b\ge 2\) 有
\[
\boxed{\ r_{a+b}\le r_a\,r_b\ }.
\]
\end{theorem}
\begin{proof}
对每个输出符号 \(c\in A_{\mathrm{out}}\)，定义线性算子
\[
\mathcal T_c^{(q)}:\mathcal H_q\longrightarrow \mathcal H_q,
\qquad
(\mathcal T_c^{(q)}p)(w):=p(B_cw).
\]
由命题 \ref{prop:pom-projective-operator-invariant-homogeneous}，
\[
\mathcal L_q|_{\mathcal H_q}=\sum_{c\in A_{\mathrm{out}}}\mathcal T_c^{(q)}.
\]
以下在各 \(\mathcal H_q\) 的单项式基下，不再区分这些算子与其矩阵表示。
\par
现在考虑乘法映射
\[
m_{a,b}:\mathcal H_a\otimes \mathcal H_b\longrightarrow \mathcal H_{a+b},
\qquad
m_{a,b}(f\otimes g):=fg.
\]
它是满射：任意次数 \(a+b\) 的单项式都可分解为一项次数 \(a\) 与一项次数 \(b\)
单项式之积。
在 \(\mathcal H_a\otimes\mathcal H_b\) 上定义
\[
\mathcal S_{a,b}:=\sum_{c\in A_{\mathrm{out}}}\mathcal T_c^{(a)}\otimes \mathcal T_c^{(b)}.
\]
则对纯张量 \(f\otimes g\) 有
\[
\begin{aligned}
m_{a,b}\!\bigl(\mathcal S_{a,b}(f\otimes g)\bigr)(w)
&=
\sum_{c\in A_{\mathrm{out}}} f(B_cw)\,g(B_cw)\\
&=
\bigl(\mathcal L_{a+b}(fg)\bigr)(w)
=
\bigl((\mathcal L_{a+b}|_{\mathcal H_{a+b}})\circ m_{a,b}\bigr)(f\otimes g)(w).
\end{aligned}
\]
故
\[
m_{a,b}\circ \mathcal S_{a,b}
=
(\mathcal L_{a+b}|_{\mathcal H_{a+b}})\circ m_{a,b}.
\]
于是 \(\mathcal H_{a+b}\) 是 \(\mathcal S_{a,b}\) 的一个商表示，从而
\[
\rho(\mathcal L_{a+b}|_{\mathcal H_{a+b}})\le \rho(\mathcal S_{a,b}).
\]
\par
固定 \(\varepsilon>0\)。
在 \(\mathcal H_q\) 的单项式基下记全 \(1\) 向量为 \(\mathbf 1_q\)，并定义
\[
x_q^{(\varepsilon)}
:=
\bigl((r_q+\varepsilon)I-(\mathcal L_q|_{\mathcal H_q})\bigr)^{-1}\mathbf 1_q
\qquad(q=a,b).
\]
由于 \(r_q+\varepsilon>\rho(\mathcal L_q|_{\mathcal H_q})\)，Neumann 级数给出
\[
x_q^{(\varepsilon)}
=
\sum_{n\ge 0}(r_q+\varepsilon)^{-n-1}
\bigl(\mathcal L_q|_{\mathcal H_q}\bigr)^n\mathbf 1_q,
\]
故 \(x_q^{(\varepsilon)}\) 的全部坐标严格为正，并满足
\[
\mathcal L_q|_{\mathcal H_q}\,x_q^{(\varepsilon)}
=(r_q+\varepsilon)x_q^{(\varepsilon)}-\mathbf 1_q
\le
(r_q+\varepsilon)x_q^{(\varepsilon)}.
\]
于是
\[
\begin{aligned}
\mathcal S_{a,b}\bigl(x_a^{(\varepsilon)}\otimes x_b^{(\varepsilon)}\bigr)
&=
\sum_c \mathcal T_c^{(a)}x_a^{(\varepsilon)}
\otimes
\mathcal T_c^{(b)}x_b^{(\varepsilon)}\\
&\le
\left(\sum_c \mathcal T_c^{(a)}x_a^{(\varepsilon)}\right)
\otimes
\left(\sum_c \mathcal T_c^{(b)}x_b^{(\varepsilon)}\right)\\
&=
\bigl(\mathcal L_a|_{\mathcal H_a}x_a^{(\varepsilon)}\bigr)
\otimes
\bigl(\mathcal L_b|_{\mathcal H_b}x_b^{(\varepsilon)}\bigr)\\
&\le
(r_a+\varepsilon)(r_b+\varepsilon)\,
x_a^{(\varepsilon)}\otimes x_b^{(\varepsilon)}.
\end{aligned}
\]
由于 \(x_a^{(\varepsilon)}\otimes x_b^{(\varepsilon)}\) 在坐标上严格为正，Collatz--Wielandt
原理给出
\[
\rho(\mathcal S_{a,b})\le (r_a+\varepsilon)(r_b+\varepsilon).
\]
合并上述两式并令 \(\varepsilon\downarrow 0\)，便得
\[
r_{a+b}
=
\rho(\mathcal L_{a+b}|_{\mathcal H_{a+b}})
\le
r_a r_b.
\]
再由定理 \ref{thm:pom-projective-operator-degenerates-to-moment-kernel}，
\(
\rho(\mathcal L_q|_{\mathcal H_q})=r_q
\)
对 \(q\ge 2\) 成立，故结论成立。证毕。
\end{proof}

\begin{corollary}[高阶谱斜率的 Fekete 极限存在]\label{cor:xi-projective-moment-fekete-slope}
沿用定理 \ref{thm:xi-projective-moment-radius-submultiplicative} 的记号，定义
\[
\tau_\ast:=\lim_{q\to\infty}\frac{1}{q}\log r_q.
\]
则该极限存在，并且
\[
\boxed{\ \tau_\ast=\inf_{q\ge 2}\frac{1}{q}\log r_q\ }.
\]
\end{corollary}
\begin{proof}
记
\[
a_q:=\log r_q
\qquad(q\ge 2),
\]
并令
\[
L:=\inf_{q\ge 2}\frac{a_q}{q}.
\]
由定理 \ref{thm:xi-projective-moment-radius-submultiplicative}，
\[
a_{m+n}\le a_m+a_n
\qquad(m,n\ge 2).
\]
因此显然有
\[
\liminf_{q\to\infty}\frac{a_q}{q}\ge L.
\]
\par
任取 \(\varepsilon>0\)，可取某个 \(q_0\ge 2\) 使
\[
\frac{a_{q_0}}{q_0}<L+\varepsilon.
\]
对任意充分大的 \(q\)，取
\[
k:=\left\lfloor\frac{q-2}{q_0}\right\rfloor,
\qquad
s:=q-kq_0.
\]
则
\[
2\le s\le q_0+1.
\]
反复应用上式的次可加性，得到
\[
a_q\le k\,a_{q_0}+C_{q_0},
\qquad
C_{q_0}:=\max_{2\le s\le q_0+1}a_s.
\]
两边除以 \(q\) 得
\[
\frac{a_q}{q}
\le
\frac{kq_0}{q}\cdot \frac{a_{q_0}}{q_0}+\frac{C_{q_0}}{q}
\le
\frac{kq_0}{q}(L+\varepsilon)+\frac{C_{q_0}}{q}.
\]
令 \(q\to\infty\)，注意 \(kq_0/q\to 1\)，便得到
\[
\limsup_{q\to\infty}\frac{a_q}{q}\le L+\varepsilon.
\]
再令 \(\varepsilon\downarrow 0\)，遂有
\[
\limsup_{q\to\infty}\frac{a_q}{q}\le L.
\]
与前述 \(\liminf\) 界合并，即得极限存在且等于 \(L\)。证毕。
\end{proof}

\begin{theorem}[原子 Witt 因子对长度同余类的贡献严格局域于 \texorpdfstring{$2\bmod m$}{2 mod m}]\label{thm:xi-atomic-witt-mod-class-localization}
沿用命题 \ref{prop:atomic-2-orbit-split-logM} 的记号。
记
\[
P(z,u)=\sum_{n\ge 1}p_n(u)\,z^n,
\qquad
P_{\mathrm{core}}(z,u)=\sum_{n\ge 1}p_n^{\mathrm{core}}(u)\,z^n,
\]
并满足精确剥离
\[
P(z,u)=u z^2+P_{\mathrm{core}}(z,u).
\]
固定模数 \(m\ge 2\)。
对每个剩余类 \(a\in \ZZ/m\ZZ\)，定义 primitive 长度同余分量
\[
P^{[a]}(z,u)
:=
\sum_{\substack{n\ge 1\\ n\equiv a\;(\mathrm{mod}\,m)}}p_n(u)\,z^n,
\]
\[
P_{\mathrm{core}}^{[a]}(z,u)
:=
\sum_{\substack{n\ge 1\\ n\equiv a\;(\mathrm{mod}\,m)}}
p_n^{\mathrm{core}}(u)\,z^n.
\]
则有严格恒等式
\[
\boxed{\;
P^{[a]}(z,u)=\delta_{a,[2]}\,u z^2+P_{\mathrm{core}}^{[a]}(z,u)\; }.
\]
等价地，若以根单位滤波定义
\[
P^{[a]}(z,u)
=
\frac1m\sum_{j=0}^{m-1}\omega_m^{-aj}P(\omega_m^j z,u),
\qquad
\omega_m:=e^{2\pi i/m},
\]
则原子 Witt 因子只在剩余类 \(a\equiv 2\pmod m\) 中留下显式项 \(u z^2\)，
对其余全部长度同余类零泄漏。
\end{theorem}
\begin{proof}
由
\[
P(z,u)=u z^2+P_{\mathrm{core}}(z,u)
\]
直接按系数比较可得
\[
p_n(u)=
\begin{cases}
u+p_2^{\mathrm{core}}(u),& n=2,\\[2mm]
p_n^{\mathrm{core}}(u),& n\neq 2.
\end{cases}
\]
故对任意剩余类 \(a\in\ZZ/m\ZZ\)，由定义立即得到
\[
\begin{aligned}
P^{[a]}(z,u)
&=
\sum_{\substack{n\ge 1\\ n\equiv a\;(\mathrm{mod}\,m)}}p_n(u)\,z^n\\
&=
\delta_{a,[2]}\,u z^2
+\sum_{\substack{n\ge 1\\ n\equiv a\;(\mathrm{mod}\,m)}}
p_n^{\mathrm{core}}(u)\,z^n\\
&=
\delta_{a,[2]}\,u z^2+P_{\mathrm{core}}^{[a]}(z,u).
\end{aligned}
\]
\par
对根单位滤波形式，只需把分解代入即可：
\[
\frac1m\sum_{j=0}^{m-1}\omega_m^{-aj}u(\omega_m^j z)^2
=
u z^2\cdot
\frac1m\sum_{j=0}^{m-1}\omega_m^{(2-a)j}
=
\delta_{a,[2]}\,u z^2.
\]
于是 Fourier 投影与按系数分组两种表述完全一致。证毕。
\end{proof}

\begin{theorem}[解析但非 RH 的实轴相确实存在]\label{thm:xi-analytic-but-not-rh-phase-realinput40}
沿用命题 \ref{prop:real-input-40-collision-pressure}、
推论 \ref{cor:real-input-40-collision-branch-radius}
与命题 \ref{prop:real-input-40-collision-rhk-window} 的记号。
记
\[
P_2(\theta):=\log\rho(e^\theta),
\qquad
R_\theta:=2.76230289346087\ldots,
\qquad
\theta_R:=\log u_R.
\]
则
\[
\boxed{\ 0<\theta_R<R_\theta\ }.
\]
特别地，开区间
\[
(\theta_R,R_\theta)\subset \RR_{>0}
\]
非空，并且对任意 \(\theta\in(\theta_R,R_\theta)\) 有：
\begin{enumerate}
  \item \(P_2(\theta)\) 仍位于 \(\theta=0\) 处 Perron 解析芽的收敛盘内部；
  \item 相应参数 \(u=e^\theta\) 已满足 \(u>u_R\)，故 kernel RH/Ramanujan 不等式失效。
\end{enumerate}
因此在真实输入 \(40\) 状态核的纯碰撞位势上，存在一个真正的
“解析但非 RH”实轴相。
\end{theorem}
\begin{proof}
由推论 \ref{cor:real-input-40-collision-branch-radius}，
\[
R_\theta\approx 2.76230289346087.
\]
由命题 \ref{prop:real-input-40-collision-rhk-window}，
\[
u_R\approx 3.59262181344,
\qquad
\theta_R=\log u_R\approx 1.27888.
\]
数值比较立即给出
\[
0<\theta_R<R_\theta.
\]
\par
若 \(\theta\in(\theta_R,R_\theta)\)，则首先 \(|\theta|<R_\theta\)，故
\(P_2(\theta)\) 仍处于 \(\theta=0\) 的解析收敛盘内。
另一方面，\(\theta>\theta_R\) 等价于 \(u=e^\theta>u_R\)，
而命题 \ref{prop:real-input-40-collision-rhk-window} 已证明
\[
\textsf{RH}(u)\Longleftrightarrow 0<u\le u_R.
\]
故当 \(u>u_R\) 时 kernel RH/Ramanujan 条件已失效。
两者合并即得结论。证毕。
\end{proof}

\begin{corollary}[总量高斯极限与同余类超平方根偏差可以严格并存]\label{cor:xi-gaussian-but-nonrh-congruence-separation}
沿用附录 \ref{app:real-input-40-zeta-u} 的输出位势记号。
令
\[
P(\theta)=\log\lambda(e^\theta),
\qquad
\lambda:=\lambda(1)=\rho(M(1)),
\]
并对每个模数 \(m\ge 2\) 与非平凡通道 \(1\le j\le m-1\) 记
\[
\rho_{m,j}:=\rho\!\bigl(M(\omega_m^j)\bigr),
\qquad
\theta_m:=\max_{1\le j\le m-1}\frac{\log \rho_{m,j}}{\log\lambda}.
\]
则在真实输入 \(40\) 状态核中同时成立：
\begin{enumerate}
  \item 无扭曲总量的输出 \(1\) 计数具有高斯极限，且其一阶 Edgeworth 修正由
  \(P^{(3)}(0)\)、\(P^{(4)}(0)\) 的闭式常数决定；
  \item 存在某个 \(m\ge 2\) 使
  \[
  \theta_m>\frac12,
  \]
  从而对应模类 Fourier 分量的衰减不满足 RH 型平方根指数。
\end{enumerate}
因此在该核上，
\[
\boxed{\;
\text{全局输出 \(1\) 计数的高斯极限}\ \not\Longrightarrow\
\text{全部同余模态的平方根消去}\; }.
\]
更强地说，在这一固定的有限状态系统内，整体输出统计呈 Gaussian--Edgeworth
行为，而某些长度/能量同余 Fourier 模态同时具有超平方根级别的偏差指数。
\end{corollary}
\begin{proof}
由命题 \ref{prop:real-input-40-output-cumulants-closed}，
\[
P^{(k)}(0)\in \QQ(\sqrt5)
\qquad(k\ge 1),
\]
并且表 \ref{tab:real_input_40_output_potential_cumulants_closed} 给出
\[
P''(0)>0.
\]
附录中的注 \ref{rem:real-input-40-output-cumulants-clt} 已说明：
输出位势的无扭曲计数方向可由同一压力函数导出中心极限定理，
而 \(P^{(3)}(0)\)、\(P^{(4)}(0)\) 则给出 Edgeworth 型的一阶修正常数。
这就是第 \(1\) 条。
\par
另一方面，命题 \ref{prop:real-input-40-output-dirichlet-nonrh} 已证明：
存在某些模数 \(m\) 使
\[
\theta_m=\max_{1\le j\le m-1}\frac{\log \rho_{m,j}}{\log\lambda}>\frac12.
\]
于是相应非平凡 Fourier 模态只以 \(O(\rho_m^n)\) 衰减，其指数严格大于
\(\lambda^{n/2}\) 所对应的平方根门槛；故第 \(2\) 条成立。
\par
最后注意：第 \(1\) 条讨论的是无扭曲总量分布在中心极限尺度上的局部形状，
而第 \(2\) 条讨论的是非平凡 Dirichlet/根单位扭曲通道的谱衰减门槛。
二者由不同的极限机制控制，并在同一核中同时成立，故得到所述严格分离。证毕。
\end{proof}

\input{subsubsec__xi-time-protocol-conclusions-part51-discrete-capacity-window6-root-center}

\endinput
